\chapter{Anytime Search: \arastar}
\label{ch:ch06}
% Function names for pseudocode are prefixed with "Ara" to stay unique
% across the book.
\SetKwFunction{AraImprovePath}{ImprovePath}
\SetKwFunction{AraMain}{ARAStar}
\SetKwFunction{AraKey}{fvalue}
\SetKwFunction{AraPublish}{Publish}
\SetKwFunction{AraPath}{PointerPath}

\begin{objectives}
\begin{itemize}
  \item Explain what an \emph{anytime} planner delivers, and why a drone that must replan within a deadline needs one rather than an optimal planner that may finish late.
  \item Trace \arastar by hand: the inflated key $\gcost + \eps\,\hcost$, the lists \Open, \Closed and \textsc{Incons}, the termination test of \AraImprovePath and the bookkeeping between two iterations.
  \item State the $\eps$-suboptimality theorem, follow its proof sketch, and compute the sharper certificate $\eps'$ from the search lists.
  \item Choose an $\eps$ schedule and predict its effect on the time to the first path, the number of published improvements and the total work, using the experiment of this chapter.
  \item Implement \arastar in \python as a generator that publishes every improved path together with its bound, and plot cost against computation time.
  \item Recognise the three classic mistakes: treating the inflated key as consistent, forgetting \textsc{Incons}, and decreasing $\eps$ too fast.
\end{itemize}
\end{objectives}

% ---------------------------------------------------------------------
\section{Why this matters for a drone swarm}
\label{sec:ch06-motivation}

A drone of the swarm in \cref{ch:ch24} is flying the path that the global
planner (\cbs or \ecbs, \cref{ch:ch09,ch:ch10}) assigned to it when an
unknown quadrotor crosses its route. The local layer (\orca or DWA,
\cref{ch:ch13,ch:ch14}) swerves, and after two seconds the drone is three
cells off its nominal path with the intruder still between it and the next
waypoint. The replanning layer must now produce a new global path from the
current cell, and it has one control cycle to do it, say $50$~ms on a flight
computer that is also running the tracker of \cref{ch:ch18} and the attitude
controller. \dstarlite (\cref{ch:ch05}) repairs the previous search quickly
when the map changed a little and the drone moved a little; after an evasive
manoeuvre both changes are large, and a repair can cost as much as a fresh
search. Plain \astar (\cref{ch:ch04}) returns an optimal path, but it returns
nothing at all until it has finished, and on a $60 \times 60$ grid with a third
of the cells blocked it expands about $1\,180$ cells before it does
(\cref{sec:ch06-experiment}). If the cycle ends first, the drone has no path.

What the drone needs is a \emph{safe path now and a better one when time
permits}. Weighted \astar with $w = 3$ delivers a path after about $215$
expansions on the same grids, on average $8\,\%$ longer than optimal, which is
a perfectly good path to start flying. The question this chapter answers is
how to keep improving that path in the remaining time \emph{without throwing
the first search away}, and how to know, at every moment, how far from optimal
the current path can be. \arastar\index{ARA*} (Anytime Repairing \astar,
Likhachev, Gordon and Thrun~\cite{likhachev2003ara}) does both: it runs
weighted \astar once with a large weight, then lowers the weight step by step
and repairs the search instead of restarting it, and after every step it
publishes the current path together with a certificate $\eps'$ such that the
path costs at most $\eps'$ times the optimum. In the study plan
(\cref{ch:appA}) it is the optional second half of Week~2, next to \dstarlite:
one planner repairs the search when the \emph{map} changes, the other when the
\emph{time budget} changes.

The chapter proceeds as follows. \Cref{sec:ch06-intuition} explains the anytime
idea and why restarting weighted \astar wastes work, \cref{sec:ch06-problem}
fixes the notation, and \cref{sec:ch06-algorithm} gives the pseudocode with a
line-by-line walkthrough. \Cref{sec:ch06-example} runs three iterations by hand
on an $8 \times 6$ grid, \cref{sec:ch06-properties} proves the suboptimality
guarantee and the sharper bound $\eps'$, and \cref{sec:ch06-schedule} studies
the choice of the $\eps$ schedule with an experiment on random grids.
\Cref{sec:ch06-variants} introduces Anytime \dstar, the combination of \arastar
with \dstarlite, \cref{sec:ch06-implementation} discusses the implementation
and its pitfalls, and \cref{sec:ch06-drone} places the algorithm in the drone
system.

% ---------------------------------------------------------------------
\section{The idea in plain words}
\label{sec:ch06-intuition}

\subsection{A path now, a better path later}

An \textbf{anytime algorithm}\index{anytime algorithm} is one that can be
interrupted at any moment and still returns a useful answer: a first, possibly
poor, solution appears quickly, the solution improves while the algorithm keeps
running, and when it is stopped the best solution found so far is handed back.
For a planner two more things are wanted. The improvements should not be
random luck but should come with a \emph{quality certificate}, a number
$\eps'$ with the meaning ``this path costs at most $\eps'$ times the optimal
cost''; and the improvements should be cheap, so that the sequence of paths
costs little more in total than a single optimal search would have. The left
panel of \cref{fig:ch06-idea} shows the resulting picture: the cost of the best
available path decreases in steps, one step per published solution, and a
deadline that falls anywhere on the time axis finds a path ready.

\begin{figure}[tb]
  \centering
  \inputfigure{ch06/idea}
  \caption{Left: an anytime planner publishes a first path early and better
  paths later; a deadline that arrives between two publications finds the last
  published path ready, together with its certificate $\eps'$. Right: the three
  lists of \arastar. States wait in \Open, are expanded once per iteration into
  \Closed, and go to \textsc{Incons} if a cheaper way to them turns up after
  they were expanded; when $\eps$ decreases, \textsc{Incons} and \Open are
  merged, their keys are recomputed and \Closed is emptied.}
  \label{fig:ch06-idea}
\end{figure}

\subsection{Weighted \astar, and why restarting it wastes work}
\label{sec:ch06-restart}

Recall from \cref{ch:ch04} that weighted \astar\index{weighted A*} is \astar
with the key $\gcost(n) + w\,\hcost(n)$ for a weight $w \ge 1$
\cite{pohl1970heuristic}. The weight makes the search greedy: cells that look
close to the goal are expanded first, the frontier stops sweeping the whole
band $\fcost \le C^*$ and dives towards the goal instead, and the search ends
after far fewer expansions. The price is bounded: with an admissible heuristic
the returned path costs at most $w \cdot C^*$, and in practice much less. On
the random grids of \cref{ch:ch04}, $w = 2$ cut the effort by a factor of
about ten while lengthening the paths by $6$--$12\,\%$.

The obvious anytime planner is therefore: run weighted \astar with $w = 3$,
publish the path, run it again with $w = 2$, publish, run it again with
$w = 1$, publish the optimal path. \Cref{sec:ch06-experiment} measures this
\textbf{restarted weighted \astar}\index{weighted A*!restarted} on random
$60 \times 60$ grids with the schedule $w = 3, 2.5, 2, 1.75, 1.5, 1.25, 1.1,
1$: it expands $3\,963$ cells in total, $3.4$ times as many as a single optimal
\astar, because every run starts from nothing and expands again most of the
cells that the previous run had already settled. Two kinds of work are
repeated. First, the $\gcost$-values are thrown away, although every one of
them is the cost of a real path and remains a valid upper bound. Second, most
of the states that the previous run expanded were expanded with a $\gcost$
that later runs never improve, so expanding them again produces nothing new.

\subsection{What can be reused}

\arastar keeps three things between runs: the $\gcost$-values with their
parent pointers, the frontier \Open, and a third list \textsc{Incons}
(\cref{fig:ch06-idea}, right). The key observation is that after a search only
the \emph{inconsistent} states matter: a state whose $\gcost$-value has been
lowered since it was last expanded, or that has never been expanded at all.
Expanding a consistent state again would relax the same edges with the same
value and change nothing. The inconsistent states are exactly the states in
\Open plus those that were improved \emph{after} they were expanded; weighted
\astar with re-opening would put the latter back into \Open at once, \arastar
records them in \textsc{Incons} and postpones their re-expansion to the next
iteration. So the next iteration starts from $\Open \cup \textsc{Incons}$ with
keys recomputed for the smaller $\eps$, and everything else, every consistent
state, keeps its value and costs nothing.

\begin{keyidea}
\arastar is weighted \astar with a decreasing weight $\eps$ that never
restarts. Each state is expanded at most once per iteration; a state improved
after its expansion goes to \textsc{Incons} instead of being re-expanded. When
$\eps$ decreases, \Open and \textsc{Incons} are merged and re-keyed, \Closed is
emptied, and the search continues from where it stopped. Every published path
costs at most $\eps' C^*$ with $\eps' = \min\bigl(\eps,\ \gcost(\gamma) /
\min_{n \in \Open \cup \textsc{Incons}} (\gcost(n) + \hcost(n))\bigr)$.
\end{keyidea}

% ---------------------------------------------------------------------
\section{Problem statement and notation}
\label{sec:ch06-problem}

We use the notation of \cref{ch:ch04}: a finite directed graph $G = (V, E)$
with edge costs $c(n, n') \ge 0$, a start $s$, a goal $\gamma$, the successor
function $\Succ(n)$, the cost $\gcost^*(n)$ of a cheapest path from $s$ to $n$,
the cost $\hcost^*(n)$ of a cheapest path from $n$ to $\gamma$, the optimal
solution cost $C^* = \gcost^*(\gamma) = \hcost^*(s)$, and a heuristic
$\hcost \colon V \to \R_{\ge 0}$ that is \textbf{consistent}: $\hcost(\gamma) =
0$ and $\hcost(n) \le c(n, n') + \hcost(n')$ on every edge. Consistency is
assumed throughout the chapter; \cref{sec:ch06-properties} shows where it is
used and \cref{exr:ch06-proof} what breaks without it. During the search,
$\gcost(n)$ is the cost of the best path from $s$ to $n$ found so far
($\infty$ if none), with a parent pointer $\mathrm{parent}(n)$ that records
the predecessor on that path. On grids the nodes are cells $(x, y)$, and the
worked example uses an 8-connected grid with move costs $1$ and $\sqrt{2}$, the
corner-cutting rule of \cref{ch:ch02} and the octile distance of
\cref{ch:ch04} as $\hcost$.

\begin{definition}[Anytime search with certificates]
\label{def:ch06-anytime}
Given a heuristic search problem and a \textbf{schedule}\index{ARA*!epsilon schedule@$\eps$ schedule} of inflation factors
$\eps_1 > \eps_2 > \dots > \eps_K = 1$, an anytime planner with certificates
outputs a sequence of solution paths $P_1, P_2, \dots, P_K$ together with
numbers $\eps'_1, \eps'_2, \dots, \eps'_K$ such that $\cost(P_k) \le \eps'_k
\cdot C^*$ for every $k$ and $\eps'_K = 1$. The planner may be interrupted
after any $P_k$; the last published pair $(P_k, \eps'_k)$ is its answer.
\end{definition}

\begin{definition}[Inflated key]
\label{def:ch06-key}
For an inflation factor $\eps \ge 1$ the \textbf{inflated
key}\index{ARA*!inflated key} of a state is
\begin{equation}
  \fcost_\eps(n) = \gcost(n) + \eps\,\hcost(n).
  \label{eq:ch06-key}
\end{equation}
For $\eps = 1$ it is the $\fcost$ of \astar; for $\eps > 1$ it is the key of
weighted \astar with $w = \eps$. In \arastar the value of $\eps$ changes between
iterations, so $\fcost_\eps$ of a state in \Open must be recomputed whenever
$\eps$ does.
\end{definition}

\begin{definition}[Consistent and inconsistent states]
\label{def:ch06-inconsistent}
Let $v(n)$ denote the value that $\gcost(n)$ had when $n$ was expanded for the
last time, and $v(n) = \infty$ if $n$ has never been expanded. A state $n$ with
$\gcost(n) < \infty$ is \textbf{consistent} if $v(n) = \gcost(n)$ and
\textbf{inconsistent}\index{inconsistent state} if $v(n) > \gcost(n)$, that is,
if a cheaper path to $n$ has been found since its last expansion, or it has
never been expanded. (Since $\gcost$ only decreases, $v(n) < \gcost(n)$ never
happens.) \arastar keeps the inconsistent states in two lists:
\Open\index{open set} holds those that may still be expanded in the current
iteration, and \textsc{Incons}\index{ARA*!INCONS list} holds those that were
already expanded in the current iteration. Together, $\Open \cup
\textsc{Incons}$ is exactly the set of inconsistent states.
\end{definition}

The word \emph{inconsistent} is used here in the sense of \lpastar and
\dstarlite (\cref{ch:ch05}), where a vertex is inconsistent when its stored
$\gcost$ disagrees with the value its predecessors would give it. It is a
different notion from a consistent \emph{heuristic}; the two meet in
\cref{sec:ch06-implementation}, where inflating a consistent heuristic is what
makes states inconsistent.

\begin{definition}[Lower bound and certificate]
\label{def:ch06-certificate}
When an iteration of \arastar ends, let
\begin{equation}
  L = \min_{n \in \Open \cup \textsc{Incons}} \bigl(\gcost(n) + \hcost(n)\bigr)
  \label{eq:ch06-lower}
\end{equation}
($L = \infty$ if both lists are empty) and define the
\textbf{certificate}\index{ARA*!suboptimality certificate}
\begin{equation}
  \eps' = \min\Bigl(\eps,\ \frac{\gcost(\gamma)}{L}\Bigr).
  \label{eq:ch06-certificate}
\end{equation}
\end{definition}

\Cref{sec:ch06-properties} shows that $L \le C^*$ unless the current path is
already optimal, so $\eps'$ is a valid bound: the published path costs at most
$\eps' C^*$. In the worked example $\eps'$ is $1.65$ after the first iteration
although $\eps = 3$, and on the random grids of \cref{sec:ch06-experiment} it
is $1.41$ for $\eps = 3$. The certificate is what the drone's executive reads;
the inflation factor is only what the search was asked for.

% ---------------------------------------------------------------------
\section{The algorithm}
\label{sec:ch06-algorithm}

\arastar consists of an inner loop \AraImprovePath, which is weighted \astar
without re-expansions (\cref{alg:ch06-improvepath}), and an outer loop
\AraMain, which runs through the schedule, does the bookkeeping between two
values of $\eps$ and publishes the paths (\cref{alg:ch06-main}). The
pseudocode follows Likhachev, Gordon and Thrun~\cite{likhachev2003ara}, with
parent pointers added so that the path can be read off.

\begin{algorithm}[htb]
\caption{\arastar, inner loop: weighted \astar without re-expansions}
\label{alg:ch06-improvepath}
\KwIn{graph $G$, goal $\gamma$, consistent heuristic $\hcost$, current $\eps$; the global tables $\gcost$, parent and the lists \Open, \Closed, \textsc{Incons}}
\KwOut{on return, $\gcost(\gamma) \le \eps\,C^*$ and the parent pointers from $\gamma$ give an $\eps$-suboptimal path}
\Fn{\AraKey{$n$}}{\KwRet $\gcost(n) + \eps\,\hcost(n)$\;}
\Fn{\AraImprovePath{}}{
  \While{$\AraKey{$\gamma$} > \min_{n \in \text{\Open}} \AraKey{$n$}$}{\label{alg:ch06-improvepath:while}
    $n \gets$ remove the state with the smallest \AraKey from \Open\label{alg:ch06-improvepath:pop}\;
    add $n$ to \Closed\label{alg:ch06-improvepath:close}\;
    \ForEach{$(n', c(n, n')) \in \Succ(n)$}{
      \If{$\gcost(n) + c(n, n') < \gcost(n')$}{\label{alg:ch06-improvepath:relax}
        $\gcost(n') \gets \gcost(n) + c(n, n')$; $\mathrm{parent}(n') \gets n$\;
        \If{$n' \notin$ \Closed}{insert $n'$ into \Open with key \AraKey{$n'$}, or update its key\label{alg:ch06-improvepath:open}}
        \Else{insert $n'$ into \textsc{Incons}\tcp*{improved after its expansion: no re-expansion now}\label{alg:ch06-improvepath:incons}}
      }
    }
  }
}
\end{algorithm}

\begin{algorithm}[htb]
\caption{\arastar, outer loop: the schedule, the bookkeeping and the certificates}
\label{alg:ch06-main}
\KwIn{graph $G$, start $s$, goal $\gamma$, consistent heuristic $\hcost$, schedule $\eps_1 > \dots > \eps_K = 1$}
\KwOut{one path with certificate per iteration; the last one is optimal}
\Fn{\AraMain{$G$, $s$, $\gamma$, $\hcost$, $(\eps_1, \dots, \eps_K)$}}{
  $\gcost(n) \gets \infty$ for all $n$; $\gcost(s) \gets 0$; parent$(s) \gets$ \KwNil\label{alg:ch06-main:init}\;
  \Open $\gets \emptyset$; \Closed $\gets \emptyset$; \textsc{Incons} $\gets \emptyset$; $\eps \gets \eps_1$\;
  insert $s$ into \Open with key \AraKey{$s$}\;
  \For{$k \gets 1$ \KwTo $K$}{
    \If{$k > 1$}{
      $\eps \gets \eps_k$\label{alg:ch06-main:decrease}\;
      move all states of \textsc{Incons} into \Open\label{alg:ch06-main:merge}\;
      recompute the key \AraKey{$n$} of every $n \in$ \Open\label{alg:ch06-main:rekey}\;
      \Closed $\gets \emptyset$\label{alg:ch06-main:clear}\;
    }
    \AraImprovePath{}\label{alg:ch06-main:improve}\;
    \If{$\gcost(\gamma) = \infty$}{\KwRet \emph{failure}\tcp*{no path exists; \Open ran empty}\label{alg:ch06-main:fail}}
    $L \gets \min_{n \in \text{\Open} \cup \textsc{Incons}} \bigl(\gcost(n) + \hcost(n)\bigr)$;\quad $\eps' \gets \min\bigl(\eps,\ \gcost(\gamma) / L\bigr)$\label{alg:ch06-main:bound}\;
    \AraPublish{\AraPath{parent, $\gamma$}, $\eps'$}\label{alg:ch06-main:publish}\;
    \If{$\eps' \le 1$}{\KwRet\tcp*{the published path is optimal; the rest of the schedule is unnecessary}\label{alg:ch06-main:stop}}
  }
}
\end{algorithm}

\paragraph{Walkthrough of \AraImprovePath.}
The loop condition on line~\ref{alg:ch06-improvepath:while} replaces the goal
test of \astar. \astar stops when it \emph{pops} the goal; \AraImprovePath
stops as soon as the goal's key is no larger than the smallest key in \Open,
which is the moment at which \astar would pop it. The difference is that the
goal is never removed from \Open: it stays there with its key, so that the next
iteration, with a smaller $\eps$, can test it again. A consequence worth
remembering is that $\gamma$ is never expanded, and once it has a finite
$\gcost$ it is always in \Open. Line~\ref{alg:ch06-improvepath:pop} pops the
best state and line~\ref{alg:ch06-improvepath:close} closes it; a closed state
is not popped again in this iteration. The relaxation on
line~\ref{alg:ch06-improvepath:relax} is that of Dijkstra's algorithm and
\astar: a strictly cheaper path to $n'$ updates $\gcost(n')$ and the parent
pointer, \emph{whatever list $n'$ is in}. The improvement is always recorded;
what depends on the list is what happens next. If $n'$ is not closed it goes
into \Open (line~\ref{alg:ch06-improvepath:open}) with a fresh key, or its key
is updated if it is there already. If $n'$ is closed
(line~\ref{alg:ch06-improvepath:incons}), \astar with re-opening would put it
back into \Open and expand it a second time; \arastar puts it into
\textsc{Incons} instead. Its successors keep their old, larger $\gcost$-values
for the rest of this iteration, which is exactly what makes the guarantee of
this iteration $\eps$-suboptimal rather than optimal, and what makes each
iteration cheap. \Cref{fig:ch06-incons} shows the situation in which
line~\ref{alg:ch06-improvepath:incons} fires in the worked example.

\begin{figure}[tb]
  \centering
  \inputfigure{ch06/incons}
  \caption{How inflation creates an inconsistent state (cells of
  \cref{ex:ch06-grid}, first iteration, $\eps = 3$). The cell $(3,3)$ can be
  reached from $(1,3)$ by two diagonal moves through $(2,2)$ (blue, cost
  $1 + 2\sqrt{2} = 3.83$) or by two straight moves through $(2,3)$ (orange,
  cost $3$). With $\eps = 3$ the diagonal neighbour has the smaller key
  ($17.41 < 18.24$), because it is closer to the goal in $\hcost$, so
  $(2,2)$ is expanded in step~3 and $(3,3)$ receives $\gcost = 3.83$ and is
  expanded in step~9. Only in step~12 is $(2,3)$ expanded and the cheaper way
  found: $(3,3)$ is closed, so its value is corrected to $3.00$ and it is moved
  to \textsc{Incons} (hatched) rather than expanded again. With $\eps = 1$ both
  neighbours would have the key $7.41$ and the cheap way would be found in
  time.}
  \label{fig:ch06-incons}
\end{figure}

\paragraph{Walkthrough of \AraMain.}
Lines~\ref{alg:ch06-main:init}--\ref{alg:ch06-main:improve} of the first pass
are weighted \astar with $w = \eps_1$ without re-opening, so the first
published path is exactly the path that weighted \astar would return, after
the same number of expansions. Between iterations four things happen
(lines~\ref{alg:ch06-main:decrease}--\ref{alg:ch06-main:clear}): $\eps$ is
decreased; the states of \textsc{Incons}, whose improvements have not yet been
propagated, are merged into \Open; the keys of all states in \Open are
recomputed for the new $\eps$, since a smaller $\eps$ changes the order of the
queue; and \Closed is emptied, so that every state may be expanded once more.
Nothing else is touched: the $\gcost$-values and parent pointers of all
states, including the goal, survive. The certificate on
line~\ref{alg:ch06-main:bound} is computed from the lists at the end of the
iteration and published with the path on line~\ref{alg:ch06-main:publish}.
Line~\ref{alg:ch06-main:stop} is an early exit: if the certificate already
says $1$, the rest of the schedule cannot improve anything, which happens
regularly in practice (\cref{sec:ch06-schedule}). If no path exists, the
first call of \AraImprovePath empties \Open, the loop condition
$\infty > \infty$ is false, and line~\ref{alg:ch06-main:fail} reports the
failure.

\paragraph{The published path.}
\AraPath follows the parent pointers back from $\gamma$ to $s$. This path may
be \emph{cheaper} than $\gcost(\gamma)$ says: when a state on it was improved
and sent to \textsc{Incons}, its parent pointer changed but the
$\gcost$-values of its descendants, including $\gcost(\gamma)$, were not
updated. The pointer path is always a valid path of cost at most
$\gcost(\gamma)$ (\cref{thm:ch06-pointer}), so publishing it can only help;
the certificate is still computed from $\gcost(\gamma)$, which is the quantity
the theorem bounds. \Cref{tab:ch06-differences} summarises how one call of
\AraImprovePath differs from the \astar of \cref{ch:ch04}.

\begin{table}[tb]
  \centering\small
  \caption{One call of \AraImprovePath compared with the graph-search \astar
  of \cref{ch:ch04}. The middle row is the defining difference; the others
  follow from it.}
  \label{tab:ch06-differences}
  \begin{tabular}{@{}lL{5.2cm}L{5.4cm}@{}}
  \toprule
  & \astar (\cref{ch:ch04}) & \AraImprovePath (\cref{alg:ch06-improvepath})\\
  \midrule
  Key of \Open & $\gcost + \hcost$ & $\gcost + \eps\,\hcost$, recomputed when $\eps$ changes\\
  Cheaper path to a closed state & re-open it and expand it again & record it, put the state into \textsc{Incons}, expand it in the next iteration\\
  Expansions per state & at most once (consistent $\hcost$) & at most once per iteration\\
  Termination & the goal is popped & $\fcost_\eps(\gamma) \le \min_{\text{\Open}} \fcost_\eps$; the goal stays in \Open\\
  Guarantee on return & optimal & $\gcost(\gamma) \le \eps\,C^*$, and $\le \eps' C^*$ with $\eps' \le \eps$\\
  State kept for the next call & nothing & $\gcost$, parents, \Open, \textsc{Incons}\\
  \bottomrule
  \end{tabular}
\end{table}

% ---------------------------------------------------------------------
\section{A worked example}
\label{sec:ch06-example}

\begin{example}[\arastar on an $8 \times 6$ grid]
\label{ex:ch06-grid}
\Cref{fig:ch06-example} shows an 8-connected $8 \times 6$ grid with ten blocked
cells, $(0,2)$, $(0,4)$, $(2,0)$, $(2,5)$, $(3,4)$, $(4,1)$, $(5,2)$, $(6,0)$,
$(6,3)$ and $(6,4)$. The start is $s = (0,3)$, the goal $\gamma = (7,2)$;
straight moves cost $1$, diagonal moves $\sqrt{2} \approx 1.41$ and are
allowed only when neither cell they cut across is blocked. The heuristic is
the octile distance, so $\hcost(s) = 7 + (\sqrt{2} - 1) \cdot 1 = 7.41$, and the
schedule is $\eps = 3, 2, 1$. Ties in the key are broken towards the smaller
$\hcost$, and among equal keys the state pushed first is popped first. The
instance is \code{example\_grid()} in \code{ch06\_arastar.py}, and every number
below is produced by its self-test. The optimal cost is $C^* = 6 + 3\sqrt{2}
= 10.24$.
\end{example}

\begin{figure}[tb]
  \centering
  \inputfigure{ch06/example}
  \caption{The three iterations of \arastar on \cref{ex:ch06-grid}. Each
  panel shows the cells expanded in that iteration (grey), the cells in \Open
  when the iteration stopped (blue), the cell moved to \textsc{Incons}
  (hatched), the published path (solid) and the previously published path
  (dashed). With $\eps = 3$ the search dives greedily along the row of the
  goal, is stopped by the wall at $x = 6$, climbs over it and publishes a path
  of cost $12.24$ after 17 expansions. With $\eps = 2$ four expansions
  propagate the correction stored in \textsc{Incons} and shorten the path to
  $11.41$. With $\eps = 1$ eleven expansions of the cheap cells that the
  greedy search had skipped find the optimal route below the wall, cost
  $10.24$.}
  \label{fig:ch06-example}
\end{figure}

\paragraph{First iteration, $\eps = 3$.}
\Cref{tab:ch06-trace} lists every expansion. The start has key $0 + 3 \cdot
7.41 = 22.24$. Its only free neighbour is $(1,3)$, whose expansion in step~2
generates five cells; the two diagonal ones, $(2,4)$ and $(2,2)$, have the
larger $\gcost$ but the smaller keys, because they are closer to the goal and
$\hcost$ counts three times. The search takes $(2,2)$, then $(3,2)$, $(4,2)$,
and the keys fall from $22.24$ to $13.41$ in five steps: this is greedy
best-first behaviour, every step reduces $\hcost$. At $(4,2)$ the wall at
$(5,2)$ blocks the way, and the key of the next expansion jumps to $15.07$
for $(4,3)$; the search climbs to $(5,4)$ in step~8, is stopped by the wall
segment $(6,3)$--$(6,4)$, and now has to fall back on cells with keys around
$17$: $(3,3)$ and $(3,1)$ in steps~9 and~10, $(4,4)$, and in step~12 the cell
$(2,3)$, which has been in \Open since step~2 with key $18.24$. Its expansion
finds the cheaper way to $(3,3)$, $\gcost = 3.00$ instead of $3.83$
(\cref{fig:ch06-incons}); $(3,3)$ is closed, so it is moved to
\textsc{Incons} and its successors, in particular $(4,3)$ with $\gcost =
4.83$, keep their values. Steps~13--17 climb over the wall along the top row
and descend the column $x = 7$ to the goal, which is generated in step~17 with
$\gcost(\gamma) = 12.24$. Its key is $12.24 + 3 \cdot 0 = 12.24$, the smallest
in \Open, and the loop stops: 17 expansions, the path of the left panel of
\cref{fig:ch06-example}, cost $8 + 3\sqrt{2} = 12.24$. The certificate: the
inconsistent states are the eight cells in \Open and $(3,3)$ in
\textsc{Incons}, and the minimum of $\gcost + \hcost$ over them is $L = 3.00 +
4.41 = 7.41$, attained by $(3,3)$; hence $\eps' = \min(3,\ 12.24 / 7.41) =
1.65$. The true ratio is $12.24 / 10.24 = 1.20$.

\begin{table}[p]
  \centering\footnotesize
  \caption{Complete trace of \cref{alg:ch06-improvepath} on
  \cref{ex:ch06-grid}, one row per expansion. Columns: the expanded cell $n$
  with its $\gcost$, $\hcost$ and key $\fcost_\eps = \gcost + \eps\,\hcost$; the
  successors whose $\gcost$ is lowered by the expansion (with the new value; a
  cell marked \textsc{I} was closed and is moved to \textsc{Incons}); and the
  smallest key in \Open after the step. An iteration stops as soon as that key
  is at least the key of the goal.}
  \label{tab:ch06-trace}
  \begin{tabular}{@{}rlrrrL{6.0cm}r@{}}
  \toprule
  Step & $n$ & $\gcost$ & $\hcost$ & $\fcost_\eps$ & Successors improved (cell$:\gcost$) & $\min_{\text{\Open}} \fcost_\eps$\\
  \midrule
  \multicolumn{7}{@{}l}{\emph{Iteration 1, $\eps = 3$; \Open holds the start only}}\\
   1 & $(0,3)$ &  0.00 & 7.41 & 22.24 & $(1,3){:}1.00$ & 20.24\\
   2 & $(1,3)$ &  1.00 & 6.41 & 20.24 & $(2,3){:}2.00$, $(1,4){:}2.00$, $(1,2){:}2.00$, $(2,4){:}2.41$, $(2,2){:}2.41$ & 17.41\\
   3 & $(2,2)$ &  2.41 & 5.00 & 17.41 & $(3,2){:}3.41$, $(2,1){:}3.41$, $(3,3){:}3.83$, $(3,1){:}3.83$, $(1,1){:}3.83$ & 15.41\\
   4 & $(3,2)$ &  3.41 & 4.00 & 15.41 & $(4,2){:}4.41$, $(4,3){:}4.83$ & 13.41\\
   5 & $(4,2)$ &  4.41 & 3.00 & 13.41 & --- & 15.07\\
   6 & $(4,3)$ &  4.83 & 3.41 & 15.07 & $(5,3){:}5.83$, $(4,4){:}5.83$, $(5,4){:}6.24$ & 13.07\\
   7 & $(5,3)$ &  5.83 & 2.41 & 13.07 & --- & 14.73\\
   8 & $(5,4)$ &  6.24 & 2.83 & 14.73 & $(5,5){:}7.24$, $(4,5){:}7.66$ & 17.07\\
   9 & $(3,3)$ &  3.83 & 4.41 & 17.07 & --- & 17.07\\
  10 & $(3,1)$ &  3.83 & 4.41 & 17.07 & $(3,0){:}4.83$ & 17.31\\
  11 & $(4,4)$ &  5.83 & 3.83 & 17.31 & $(4,5){:}6.83$ & 18.24\\
  12 & $(2,3)$ &  2.00 & 5.41 & 18.24 & $(3,3){:}3.00$ \textsc{I} & 18.73\\
  13 & $(5,5)$ &  7.24 & 3.83 & 18.73 & $(6,5){:}8.24$ & 18.49\\
  14 & $(6,5)$ &  8.24 & 3.41 & 18.49 & $(7,5){:}9.24$ & 18.24\\
  15 & $(7,5)$ &  9.24 & 3.00 & 18.24 & $(7,4){:}10.24$ & 16.24\\
  16 & $(7,4)$ & 10.24 & 2.00 & 16.24 & $(7,3){:}11.24$ & 14.24\\
  17 & $(7,3)$ & 11.24 & 1.00 & 14.24 & $(7,2){:}12.24$ & 12.24\\
  \multicolumn{7}{@{}l}{\quad stop: $\fcost_3(\gamma) = 12.24 \le 12.24$; $L = 7.41$ from $(3,3)$, $\eps' = 1.65$}\\
  \midrule
  \multicolumn{7}{@{}l}{\emph{Iteration 2, $\eps = 2$; \Open holds the eight open cells and $(3,3)$ from \textsc{Incons}, re-keyed}}\\
   1 & $(3,3)$ &  3.00 & 4.41 & 11.83 & $(4,3){:}4.00$ & 10.83\\
   2 & $(4,3)$ &  4.00 & 3.41 & 10.83 & $(5,3){:}5.00$, $(4,4){:}5.00$, $(5,4){:}5.41$ & 9.83\\
   3 & $(5,3)$ &  5.00 & 2.41 &  9.83 & --- & 11.07\\
   4 & $(5,4)$ &  5.41 & 2.83 & 11.07 & $(5,5){:}6.41$ & 12.24\\
  \multicolumn{7}{@{}l}{\quad stop: $\fcost_2(\gamma) = 12.24 \le 12.24$; $L = 8.00$ from $(1,2)$, $\eps' = 1.53$; pointer path costs $11.41$}\\
  \midrule
  \multicolumn{7}{@{}l}{\emph{Iteration 3, $\eps = 1$; \Open holds ten cells, re-keyed}}\\
   1 & $(1,2)$ &  2.00 & 6.00 &  8.00 & $(1,1){:}3.00$ & 8.24\\
   2 & $(2,4)$ &  2.41 & 5.83 &  8.24 & --- & 8.83\\
   3 & $(4,4)$ &  5.00 & 3.83 &  8.83 & $(4,5){:}6.00$ & 8.83\\
   4 & $(2,1)$ &  3.41 & 5.41 &  8.83 & --- & 8.83\\
   5 & $(1,4)$ &  2.00 & 6.83 &  8.83 & $(1,5){:}3.00$ & 9.41\\
   6 & $(1,1)$ &  3.00 & 6.41 &  9.41 & $(0,1){:}4.00$, $(1,0){:}4.00$, $(0,0){:}4.41$ & 9.66\\
   7 & $(3,0)$ &  4.83 & 4.83 &  9.66 & $(4,0){:}5.83$ & 9.66\\
   8 & $(4,0)$ &  5.83 & 3.83 &  9.66 & $(5,0){:}6.83$ & 9.66\\
   9 & $(5,0)$ &  6.83 & 2.83 &  9.66 & $(5,1){:}7.83$ & 10.24\\
  10 & $(5,1)$ &  7.83 & 2.41 & 10.24 & $(6,1){:}8.83$ & 10.24\\
  11 & $(6,1)$ &  8.83 & 1.41 & 10.24 & $(7,1){:}9.83$, $(6,2){:}9.83$, $(7,2){:}10.24$ & 10.24\\
  \multicolumn{7}{@{}l}{\quad stop: $\fcost_1(\gamma) = 10.24 \le 10.24$; $L = 10.24$, $\eps' = 1$: optimal}\\
  \bottomrule
  \end{tabular}
\end{table}

\paragraph{Second iteration, $\eps = 2$.}
The nine inconsistent states are re-keyed with $\eps = 2$. The key of the goal
is still $12.24$; the smallest key belongs to $(3,3)$ from \textsc{Incons},
$3.00 + 2 \cdot 4.41 = 11.83$, and it is expanded first. This is the postponed
re-expansion: it lowers $(4,3)$ from $4.83$ to $4.00$, and the next three
expansions carry the correction on to $(5,3)$, $(4,4)$, $(5,4)$ and $(5,5)$,
each of which is in \Open again because \Closed was emptied. After step~4 the
smallest key in \Open is $12.24$, the goal's own key: the cells at the top of
the wall, $(5,5)$ and $(6,5)$ onwards, are not re-expanded because their keys
$6.41 + 2 \cdot 3.83 = 14.07$ and higher exceed it. So $\gcost(\gamma)$ stays
at $12.24$, but the parent pointers now run $\gamma \to (7,3) \to \dots \to
(5,5) \to (5,4) \to (4,3) \to (3,3) \to (2,3) \to (1,3) \to s$, and the pointer
path costs $10 + \sqrt{2} = 11.41$: the published path is cheaper than its own
$\gcost$-value, the effect described in \cref{sec:ch06-algorithm}. Four
expansions, against 20 for weighted \astar with $w = 2$ started from scratch,
and the certificate improves to $\eps' = \min(2,\ 12.24 / 8.00) = 1.53$, with
$L = 8.00$ now attained by $(1,2)$.

\paragraph{Third iteration, $\eps = 1$.}
With $\eps = 1$ the keys are the plain $\fcost$-values of \astar, and the cells
that the greedy first iteration skipped, $(1,2)$, $(2,4)$, $(2,1)$, $(1,4)$
and $(1,1)$ with $\fcost$ between $8$ and $9.41$, are expanded first, most of
them without improving anything: this is the price of the guarantee, \astar
must rule them out. From step~7 on the search follows the bottom corridor
$(3,0), (4,0), (5,0), (5,1), (6,1)$, whose cells had been generated in the
first iteration with their optimal $\gcost$-values but never expanded, and
step~11 generates the goal with $\gcost(\gamma) = 10.24$. The smallest key in
\Open is now $10.24$, the loop stops, and $L = 10.24$ gives $\eps' = 1$: the
path is optimal. Eleven expansions, 32 in total for the three iterations.
Plain \astar with the same tie-breaking needs 22 expansions for this instance,
and restarting weighted \astar for each $\eps$ would need $17 + 20 + 22 = 59$.
\Cref{tab:ch06-iterations} collects the numbers.

\begin{table}[tb]
  \centering\small
  \caption{The three iterations of \cref{ex:ch06-grid}: expansions, the
  sizes of the lists when the iteration stops, the value $\gcost(\gamma)$ and
  the cost of the pointer path, the lower bound $L$, the certificate $\eps'$
  and the true ratio to $C^* = 10.24$. In iteration~2 the pointer path is
  cheaper than $\gcost(\gamma)$.}
  \label{tab:ch06-iterations}
  \begin{tabular}{@{}ccrrccrrrrr@{}}
  \toprule
  $k$ & $\eps_k$ & Exp. & Total & $\abs{\text{\Open}}$ & \textsc{Incons} & $\gcost(\gamma)$ & Path cost & $L$ & $\eps'_k$ & Cost$/C^*$\\
  \midrule
  1 & 3 & 17 & 17 &  8 & $\{(3,3)\}$ & 12.24 & 12.24 &  7.41 & 1.65 & 1.20\\
  2 & 2 &  4 & 21 & 10 & $\emptyset$ & 12.24 & 11.41 &  8.00 & 1.53 & 1.11\\
  3 & 1 & 11 & 32 &  9 & $\emptyset$ & 10.24 & 10.24 & 10.24 & 1.00 & 1.00\\
  \bottomrule
  \end{tabular}
